\chapter{FPGA 实时性能与资源效率}
\label{chap:verification}

本章围绕两个问题评价 FPGA 实现：一是能否在持续视频输入与显示的同时，
按期完成完整矫正帧；二是为达到这一吞吐，消耗了多少计算、存储与互连资源。
分析采用原始 RTL 仿真日志、PDS 综合和布局布线报告，以及
\path{chapter6/references/} 中的性能说明与审计记录。
其中，1080p 完整链路仿真给出了约 30 fps 的新结果发布率，
最慢观测渲染时间为 33.14614 ms；128$\times$32 分块实现则将物理 DRM
占用从历史实现的 122 个位置降至 104 个。
这两项结果分别反映已测试条件下的处理能力和存储效率，其适用版本与测量条件如下。

\section{分析配置与数据来源}
\subsection{工作配置及评价指标}
系统基于紫光同创 PGL50H-FBG484-6 FPGA。默认数据通路接收
$1920\times1080$ RGB888 视频，依次执行帧自适应 Gamma、
RGB565 量化、DDR 源帧存储、姿态与镜头畸变联合逆投影、最近邻重采样和结果写回，
再将 RGB565 展开为 RGB888 输出。输入与输出均使用 148.5 MHz 像素时钟和
60 Hz 扫描时序，矫正器的目标是每秒生成 30 个来自不同源帧的完整结果。

\begin{table}[htbp]
\centering\small
\caption{实时性能分析采用的默认配置}\label{tab:rt-config}
\begin{tabularx}{\textwidth}{p{4.4cm}Y}
\toprule
配置项 & 参数及含义 \\
\midrule
有效图像与总光栅 & $1920\times1080$；$2200\times1125$ \\
像素时钟与扫描频率 & 148.5 MHz；60 Hz \\
DDR 用户接口与渲染时钟 & 256 bit；100 MHz，每 beat 为 32 bytes \\
DDR 像素格式 & RGB565，每 beat 承载 16 个像素 \\
输出分块与片内块槽 & $128\times32$；两个块槽交替使用 \\
单槽源区域缓存 & $2048\times256$ bit，即 64 KiB \\
片外帧槽 & 两个源帧槽、两个结果帧槽 \\
几何计算与采样 & 逐像素定点逆投影；最近邻 \\
姿态配置 & yaw、pitch 各约 $\pm15^\circ$；roll 未参与变换 \\
结果交接 & 最后一笔写数据传输完成后提交，在显示帧边界切换 \\
\bottomrule
\end{tabularx}
\end{table}

本章区分扫描频率、新结果帧率、单帧渲染时间和端到端延迟。
扫描频率描述接口发送图像的节拍；新结果帧率描述不同源帧的处理产出；
渲染时间描述矫正器占用一帧任务的时间；端到端延迟还包含采集、调度与显示等待。
资源效率则分别用综合逻辑用量、最终物理占用、逻辑缓存容量和版本间变化评价。

\subsection{报告版本与证据边界}
\path{../Pgl50h/prj_tasks/} 中的报告头部标明 PDS2025.2 build 211867，
生成日期为 2026 年 10 月 21 日，其综合用量为 17,294 LUT、25,891 FF、
121 个 DRM18K 等效容量。它们属于较早的高存储占用实现，
不能作为 128$\times$32 默认版本的资源表。
表~\ref{tab:rt-evidence} 将本章使用的数据按生成日期分开。

\begin{table}[htbp]
\centering\small
\caption{性能与资源数据的版本划分}\label{tab:rt-evidence}
\begin{tabularx}{\textwidth}{p{2.1cm}>{\raggedright\arraybackslash}p{3.2cm}Y}
\toprule
证据日期 & 来源与工具 & 本章用途 \\
\midrule
2026-10-21 & 工程 \path{prj_tasks/}；PDS2025.2 & 直接读取综合、物理资源与双角时序，作为早期实现对照。\\
2026-10-22 & 完整链路仿真及配套实现归档；PDS2025.2 & 分析 $128\times32$ 的新帧率、渲染耗时、FIFO 峰值及共享写缓冲收益。\\
2026-10-24 & 本章参考说明与数值审计；PDS2025.2 & 引用复核版本的资源和时序结果，独立于 10 月 22 日的吞吐数据。\\
2026-10-25 & 本章 JSON 审计；PDS2022.2-SP3 & 说明跨工具重编译的时序失败，以及未采用修复候选的验证状态。\\
\bottomrule
\end{tabularx}
\end{table}

10 月 25 日审计将 10 月 22 日完整 1080p 测试明确标记为历史结果，
并列出了后续变化的帧缓存、渲染器、FIFO 与测试文件。
因此，30.000030 fps 是对应归档快照的结果，尚不能直接写成后续源码的重新测试值。
审计还记录了用户反馈的仓库版本板测状态，但没有核验实际烧录位流哈希，
也没有提供板级帧率或端到端延迟数值；本章据此保留功能反馈与定量测量之间的区别。
\begingroup\interlinepenalty=10000
\sourcefile{\path{chapter6/references/} 下的《02\_实时性能与资源效率.md》；
\path{chapter6/references/realtime_resource_audit_20261025.json}；
\path{../Pgl50h/prj_tasks/syn_1/synthesize/hdmi_ddr3_top.snr}}
\endgroup

\section{实时处理预算与流水重叠}
\subsection{帧率目标与单帧预算}
一帧有效图像包含
\begin{equation}
N=1920\times1080=2,073,600\ \text{pixels}.
\end{equation}
输入和显示扫描频率为
\begin{equation}
f_{\rm scan}=\frac{148.5\times10^6}{2200\times1125}
=60\ \text{Hz},\qquad T_{\rm scan}=16.66667\ \text{ms}.
\end{equation}
在此扫描时序下，有效输入像素平均速率为 124.416 Mpixel/s。
30 fps 矫正目标对应的有效处理需求与 100 MHz 域单帧预算为
\begin{align}
P_{30}&=N\times30=62.208\ \text{Mpixel/s},\\
C_{30}&=\frac{100\times10^6}{30}=3,333,333.33\ \text{cycles},\qquad
T_{30}=33.333333\ \text{ms}.
\end{align}

100 MHz 下每拍完成一个像素的单流水线，理想吞吐上限为 100 Mpixel/s，
折合完整 1080p 图像约 48.23 fps。因而默认单像素通路需要依靠流水重叠达到
30 fps；若要求完整 1080p60，即使忽略全部调度与访存开销，
100 MHz、每拍一个像素的输出通路也不足以满足 124.416 Mpixel/s 的需求。
60 Hz HDMI 可以重复显示同一完整结果帧，其连续输出并不增加矫正器的新帧产出。

\subsection{块内流水与块间并行}
$128\times32$ 分块每块最多 4096 个输出像素，一帧共有
\begin{equation}
N_{\rm tile}=\left\lceil\frac{1920}{128}\right\rceil
\left\lceil\frac{1080}{32}\right\rceil=15\times34=510\ \text{blocks}.
\end{equation}
末排高度为 24，按实际有效高度处理。每块逐像素生成逆投影坐标，
再按源包围矩形预取数据，随后采样并写回；块内像素计算使用流水线，
两个块槽则使相邻块的不同阶段重叠。

\begin{figure}[htbp]
\centering
\begin{tikzpicture}[x=1cm,y=1cm,font=\small,
  phase/.style={draw=TemplateBlue,rounded corners=2pt,minimum height=0.65cm,
    text width=3.1cm,align=center}]
\draw[-{Latex}] (1.7,0.55)--(12.9,0.55) node[above left]{处理推进};
\node[left] at (1.6,0) {块 $k$};
\node[left] at (1.6,-0.9) {块 $k+1$};
\node[left] at (1.6,-1.8) {块 $k+2$};
\node[phase,fill=TemplateBlue!8] at (3.5,0) {投影与源区域预取};
\node[phase,fill=TemplateBlue!22] at (7.0,0) {采样与结果写回};
\node[phase,fill=TemplateBlue!8] at (7.0,-0.9) {投影与源区域预取};
\node[phase,fill=TemplateBlue!22] at (10.5,-0.9) {采样与结果写回};
\node[phase,fill=TemplateBlue!8] at (10.5,-1.8) {投影与源区域预取};
\end{tikzpicture}
\caption{双块槽的阶段重叠示意。图中长度不表示实测时间；源读取与写回仍共享 DDR 服务能力。}
\label{fig:rt-overlap}
\end{figure}

若将块的投影及源加载耗时记作 $C_{\rm front}$，采样及结果写回耗时记作
$C_{\rm back}$，在资源互不争用的理想条件下，块间启动间隔可接近
$\max(C_{\rm front},C_{\rm back})$，而非两者之和。
实际系统中读写竞争、源矩形大小、块槽释放和流水排空会延长间隔，
因此图~\ref{fig:rt-overlap} 描述的是可重叠关系，不能用理想阶段时间替代整帧测量。

帧级并行进一步使采集下一幅源图、渲染已提交源帧和显示已完成结果同时进行。
两个源帧槽保护正在渲染的图像，两个结果帧槽隔离显示与写回。
渲染器只有在最后一笔有效写传输完成后才提交结果，显示侧在帧边界完成交接。
这使重复显示和输入跳帧成为吞吐调节方式，同时保证正在显示的图像不会被局部覆盖。
\begingroup\interlinepenalty=10000
\sourcefile{\path{../Pgl50h/src/video/pose/fram_buf_pose.v}；
\path{../Pgl50h/src/video/pose/pose_renderer_tiled.v}}
\endgroup

\section{完整链路吞吐与时间余量}
\subsection{仿真条件及观测结果}
10 月 22 日归档仿真使用 ModelSim SE-64 10.6e，实例化实际帧缓存、
异步 FIFO、渲染器与仲裁器，同时运行 Gamma、RGB565 转换、采集和显示。
DDR 行为模型设置 32 拍命令延迟、25\% 随机数据停顿、
读写共享每拍最多一个 256 bit 数据 beat，并每 780 拍暂停 35 拍模拟刷新。
该模型用于施加可重复的访存压力，未包含厂商 DDR PHY 和物理 UART。

\begin{table}[htbp]
\centering\small
\caption{10 月 22 日完整 1080p 并发 RTL 仿真结果}\label{tab:simulation}
\begin{tabularx}{\textwidth}{p{5.6cm}Y}
\toprule
指标 & 观测结果 \\
\midrule
输入帧 / 不同源帧结果发布数 & 28 / 13 \\
稳态新结果发布率 & 30.000030 fps \\
相邻新结果发布间隔 & 全部为 3,333,330 拍，即 33.33330 ms \\
最慢单帧渲染 & 3,314,614 拍，即 33.14614 ms \\
13 次渲染的最小值 / 平均值 & 28.70689 ms / 31.75641 ms \\
采集 FIFO 峰值 & 15/512 个 256 bit 字，即 480 bytes \\
显示像素比较次数 & 51,840,000 \\
并行与竞争记录 & 重叠 20,080,265 拍；共享通道竞争计数 396,039 \\
测试覆盖 & 11 组姿态激励、全部 8 种控制掩码 \\
日志结果 & PASS，0 errors，4 warnings \\
\bottomrule
\end{tabularx}
\end{table}

新帧率按首末完整新结果发布时刻计算。设发布数为 $M$、对应 100 MHz 计数为
$C_1$ 与 $C_M$，则
\begin{equation}
f_{\rm new}=\frac{(M-1)f_{\rm DDR}}{C_M-C_1}
=\frac{12\times100\times10^6}{45,000,053-5,000,093}
=30.000030\ \text{fps}.
\end{equation}
该值略高于 30 来自仿真时钟离散，实际意义为达到约 30 fps 的发布节拍。
28 个输入帧与 13 个结果包含启动、稳态与排空阶段，不能用两者直接相除定义稳态丢帧率。
日志中的源帧编号也显示系统选择部分源帧处理，显示端按 60 Hz 时序重复读出完整结果。

表中最小值与平均值由 13 条 \texttt{RENDER} 记录重新统计，包含启动和不同控制状态，
只说明该次运行的样本分布。姿态掩码与功能掩码覆盖并不等于两者所有组合均完成验证；
部分姿态激励下几何功能处于关闭状态。四条 warning 均与测试平台 Gamma
统计元数据端口连接有关，因此本次像素检查也不能代替完整遥测链路的验证。
\begingroup\interlinepenalty=10000
\sourcefile{\path{../Pgl50h/verification/pds_pc_control_write_ram_20261022/simulation/result.log}}
\endgroup

\subsection{整帧效率与裕量解释}
相对 30 fps 的单帧预算，最慢观测渲染的时间余量为
\begin{align}
\Delta C&=C_{30}-3,314,614=18,719.33\ \text{cycles},\\
\Delta T&=T_{30}-33.14614=0.187193\ \text{ms},\\
\eta_{\rm margin}&=\frac{\Delta T}{T_{30}}\times100\%=0.56158\%.
\end{align}
该帧平均每个输出像素消耗的周期数及有效吞吐为
\begin{equation}
c_{\rm pixel}=\frac{3,314,614}{2,073,600}=1.59848,
\qquad
P_{\rm observed}=\frac{2,073,600}{33.14614\times10^{-3}}
\approx62.5593\ \text{Mpixel/s}.
\end{equation}
这一整帧指标已经包含投影、源读取、采样、写回与等待，
不能将其中多出的 0.59848 拍/像素全部归因于 DDR 停顿。
相对单像素每拍的 100 Mpixel/s 理想上限，观测有效吞吐约为 62.56\%；
它反映整个调度与数据通路的工作效率，不是算术单元利用率。

0.56\% 的帧预算余量说明测试已达到目标，但对附加等待较敏感。
例如，增加 20,000 个不能被重叠隐藏的 100 MHz 周期将增加 0.2 ms，
足以耗尽这一余量。若结果错过显示帧边界，新图像可能推迟到下一次扫描交接，
从而使输出更新时间产生离散跳变。因此，平均帧率之外，还需要检查最大发布间隔、
最大渲染时间和显示欠载事件，不能仅凭平均吞吐判定所有场景都能按期完成。

FIFO 峰值 15/512 对应深度占用 2.93\%，说明已覆盖负载下采集侧缓冲充足。
但它只约束该次仿真的队列水位，不能推出真实 DDR 的最长无服务时间，
也不能推出渲染阶段具有相同幅度的时间余量。已有结果是有限姿态和特定压力模型下的
观测最大值，尚未构成全部姿态、全部存储调度下的最坏执行时间证明。

\section{处理延迟与显示响应}
\subsection{端到端延迟的事件定义}
本系统先提交完整源帧，再渲染并交接完整结果，因而单帧渲染时间只是总延迟的一部分。
以同一被接受源帧的输入 SOF，到其对应结果首个有效输出像素为起止事件，定义
\begin{equation}
T_{\rm e2e}=T_{\rm capture}+T_{\rm queue}+T_{\rm render}
+T_{\rm handoff}+T_{\rm active}.
\label{eq:rt-latency}
\end{equation}

\begin{table}[htbp]
\centering\small
\caption{端到端延迟分项及其统计事件}\label{tab:rt-latency}
\begin{tabularx}{\textwidth}{p{2.5cm}Y}
\toprule
分项 & 起止事件与影响因素 \\
\midrule
$T_{\rm capture}$ & 接受源帧 SOF 至该帧末笔采集写完成；包括整幅有效图像到达与输入侧排空。\\
$T_{\rm queue}$ & 源帧提交至渲染启动；受渲染器、帧槽及结果发布状态影响。\\
$T_{\rm render}$ & 渲染启动至末笔结果写完成；归档观测最大值为 33.14614 ms。\\
$T_{\rm handoff}$ & 结果提交至显示侧完成帧边界交接；包括旧读事务排空与跨域握手。\\
$T_{\rm active}$ & 完成交接至该结果首个有效像素输出；取决于消隐区及预取状态。\\
\bottomrule
\end{tabularx}
\end{table}

内部写完成表示 DDR IP 用户接口正确接受了最后一个数据 beat，
该接口没有独立写响应 B 通道；它不是对 DDR 芯片物理写入时刻的直接测量。
延迟统计应始终沿用同一接口事件定义。
若改用输入/输出 SOF 或指定位置像素作为端点，式~\eqref{eq:rt-latency}
中的扫描等待项也应相应调整，避免重复计算消隐时间。

输入寄存、Gamma 和 RGB565 转换共四拍，在 148.5 MHz 下约为 26.94 ns，
远小于毫秒级整帧缓存和渲染耗时。缩短少量输入寄存级对总响应改善有限；
减少整帧等待、提升渲染速度和避免错过显示交接窗口，才会显著影响输出响应。
在能够赶上下一次交接且没有附加阻塞的条件下，帧边界等待通常小于一个
16.67 ms 扫描周期，但现有资料不足以将该条件推广为总延迟上界。

\subsection{已有数据能够支持的延迟结论}
现有完整链路日志记录渲染和发布事件，未给出同源输入 SOF 与首个输出有效像素的
完整配对统计，故本章不将渲染耗时填作板级端到端延迟。
测量时应为接受的源帧赋予递增身份，在统一参考时钟下记录源帧接受、提交、
渲染启动、末写完成、显示槽切换及输出有效事件，匹配同一身份后计算各分项。
未被系统接受的输入帧应计入跳帧统计，不能作为延迟样本。

板级测量还应明确是否包含相机曝光、HDMI 桥、采集卡和显示设备延迟，
预热后连续记录不少于 60 s，报告样本数、平均值、P95、最大值及异常事件。
当前审计中的板级帧率与端到端延迟字段均为空，尚无可引用的定量实测值。

\section{DDR 访问效率与带宽约束}
\subsection{有效载荷与带宽预算}
一帧 RGB565 图像的有效载荷及 beat 数为
\begin{equation}
S_f=2N=4,147,200\ \text{bytes}\approx3.955\ \text{MiB},\qquad
B_f=\frac{S_f}{32}=129,600\ \text{beats}.
\end{equation}
256 bit、100 MHz 用户数据通道的理想载荷上限为
\begin{equation}
B_{\rm peak}=32\times100\times10^6=3.2\ \text{GB/s}.
\end{equation}
该值按一次数据 beat 计算，未扣除命令、读写切换、刷新、仲裁和背压，
也不表示真实板卡的持续有效带宽。以下带宽采用十进制 MB/s，存储容量采用二进制 MiB。

设每个矫正帧实际读取 $R$ 个源区域 beat，目标渲染率为 30 fps，则
\begin{table}[htbp]
\centering\small
\caption{默认路径的 DDR 有效数据预算}\label{tab:rt-bandwidth}
\begin{tabularx}{\textwidth}{p{3cm}Y>{\raggedright\arraybackslash}p{3.8cm}}
\toprule
数据流 & 预算条件 & 有效载荷带宽 \\
\midrule
源帧采集写入 & 按 60 fps 全部接受计算 & 248.832 MB/s \\
结果帧显示读取 & 每秒扫描读取 60 个完整帧 & 248.832 MB/s \\
矫正结果写回 & 每秒写回 30 个完整帧 & 124.416 MB/s \\
渲染源区域读取 & 每帧读取 $R$ beats & $0.000960R$ MB/s \\
合计 & 前三项预算加姿态相关读取 & $622.080+0.000960R$ MB/s \\
\bottomrule
\end{tabularx}
\end{table}

若仅作理想示例，假设每个源像素每帧恰好读取一次，则 $R=B_f$，
总有效载荷为 746.496 MB/s，约占接口峰值的 23.33\%。
这一比例不能解释为实测 DDR 利用率：当前分块加载可能重复读取相邻源区域，
无效块又可能不读源图，实际采集也可能跳过输入帧。
如果用 $\alpha=R/B_f$ 表示相对一次整帧读取的源载荷系数，则预算可写为
\begin{equation}
B_{\rm total}=622.080+124.416\alpha\quad\text{MB/s}.
\end{equation}
当前 $128\times32$ 完整系统日志未给出逐姿态的 $R$，
因此保留这一参数，而不将历史 $128\times64$ 的源读取数填入当前配置。

\subsection{分块、突发长度与命令开销}
渲染器根据实际采样坐标求源包围矩形，并沿水平方向对齐到 16 像素边界。
设有效坐标极值为 $x_{\min},x_{\max},y_{\min},y_{\max}$，则单块源数据需求为
\begin{equation}
B_s=\left(\left\lfloor\frac{x_{\max}}{16}\right\rfloor
-\left\lfloor\frac{x_{\min}}{16}\right\rfloor+1\right)
(y_{\max}-y_{\min}+1),\qquad B_s\le2048.
\end{equation}
矩形加载将不规则像素访问转换为连续 burst，代价是读取矩形内未被采样的像素、
横向对齐填充，以及相邻块重叠区域。姿态和镜头参数改变矩形的形状，
因此源读取量和处理耗时均随场景变化。

内部 burst 长度为 1～16 beat。完整 $128$ 像素输出行包含 8 beat，
最大 $128\times32$ 块包含 256 beat；按当前逐行写回组织，
全帧结果约需 $15\times1080=16,200$ 笔 8-beat 写事务。
将块高从 64 减为 32，并不直接改变每行 8 beat 的打包方式，
但使块级投影准备、容量检查和槽切换次数从 255 增至 510，
也可能增加源矩形边界的重复读取。

以不重叠命令与数据的孤立事务为例，忽略随机停顿时，32 拍命令等待与 $L$ 拍载荷
对应的效率为 $L/(32+L)$：8-beat 约为 20\%，16-beat 约为 33.33\%。
这只是说明短突发对等待敏感的简化模型；实际完整系统存在多阶段重叠，
不能把该比例直接乘以接口峰值作为实测持续带宽。
模型中的 25\% 数据停顿与刷新还会进一步减少可服务周期，
两种停顿是否重合也会影响结果，故不能将各比例简单相加得到总损失。

\subsection{带宽充足与按期完成的区别}
显示端必须按有效像素节拍持续供数，采集端必须及时接收输入；
两者对服务间隔敏感，而渲染端还要在 33.33 ms 内完成一帧。
即使平均数据量低于 3.2 GB/s，集中出现的命令等待、刷新或仲裁也可能造成
FIFO 欠载或错过发布窗口。归档中的共享竞争与刷新暂停计数证明确实存在等待，
但这些计数会与其他流水活动重叠，不能全部线性累加为单帧延迟。

因此，DDR 效率应结合有效 beat 数、命令数、平均 burst 长度、各客户端等待时间、
FIFO 最高/最低水位以及帧完成期限评价。当前已有完整链路压力测试与采集 FIFO 峰值，
能够支持已覆盖条件下的吞吐结论；板级持续带宽及各姿态访问放大仍需对应事件计数。

\section{资源占用与存储组织效率}
\subsection{综合逻辑与最终物理资源}
表~\ref{tab:rt-synthesis} 列出三个 PDS2025.2 版本的综合结果。
前两列直接读取原始 SNR，10 月 24 日列采用本章参考资料中的数值。
APM 是 PDS 的 DSP 类算术资源名称，DRM18K 等效量按存储块容量折算；
分布式 RAM LUT 已包含在总 LUT 中，不能再次相加。

\begin{table}[htbp]
\centering\small
\caption{不同版本的综合资源用量}\label{tab:rt-synthesis}
\begin{tabularx}{\textwidth}{Yrrrr}
\toprule
资源 & 10 月 21 日 & 10 月 22 日 & 10 月 24 日 & 器件容量 \\
\midrule
LUT & 17,294 & 14,571 & 14,565 & 42,800 \\
FF & 25,891 & 15,006 & 15,010 & 64,200 \\
APM & 45 & 45 & 45 & 84 \\
DRM18K 等效量 & 121 & 103 & 103 & 134 \\
分布式 RAM LUT & 82 & 594 & 594 & 17,000 \\
锁存器 & 0 & 0 & 0 & --- \\
\bottomrule
\end{tabularx}
\end{table}

物理报告进一步包含常量、旁路、实际存储位置以及映射后的寄存器组织。
表~\ref{tab:rt-physical} 的 FF 均使用 CLMA 与 CLMS 内 FF 之和，不包含 I/O 寄存器。
10 月 21 日的 LUT 为 $11,914+7,438=19,352$，
FF 为 $15,911+9,845=25,756$；10 月 22 日对应为
$10,181+6,697=16,878$ 和 $9,312+5,559=14,871$。

\begin{table}[htbp]
\centering\small
\caption{布局布线后的物理资源及 10 月 24 日版本余量}\label{tab:rt-physical}
\begin{tabularx}{\textwidth}{Yrrrrr}
\toprule
资源 & 10 月 21 日 & 10 月 22 日 & 10 月 24 日 & 占用率 & 剩余 \\
\midrule
LUT & 19,352 & 16,878 & 16,897 & 39.48\% & 25,903 \\
CLM FF & 25,756 & 14,871 & 14,875 & 23.17\% & 49,325 \\
APM & 45 & 45 & 45 & 53.57\% & 39 \\
DRM 位置 & 122 & 104 & 104 & 77.61\% & 30 \\
\bottomrule
\end{tabularx}
\end{table}
占用率与剩余量对应 10 月 24 日列，按各类器件容量重新计算。

10 月 22 日与 24 日的 DRM 均由 102 个 DRM18K 和 2 个 DRM9K 构成。
按容量为 $102+2/2=103$ 个 DRM18K 等效量，利用率 76.87\%；
按位置则为 $102+2=104$ 个，利用率 77.61\%。
9K 模式占据物理位置，剩余容量也不能任意拼接，因此两种统计不能互换。
原始物理表常使用向上取整的整数百分比，例如 $122/134$ 显示为 92\%；
本章统一以计数重新计算为 91.04\%，避免将显示取整误认为资源变化。

10 月 24 日物理 LUT 比综合 LUT 多 2332 个，FF 则少 135 个。
这反映了映射、常量及旁路计数和寄存器统计边界的差异，
不宜将综合与物理列的差额直接解释为新增算法开销。
从主要数据资源看，DRM 是最紧张的一类，APM 次之，LUT 和 FF 相对宽裕。
此外，10 月 22 日 ROR 中 BKCL 已使用 4/4；
后续功能扩展还受专用时钟资源、布局位置和布线能力约束，不能仅按空余 LUT 推断可扩展规模。
\begingroup\interlinepenalty=10000
\sourcefile{\path{../Pgl50h/prj_tasks/pnr_1/place_route/hdmi_ddr3_top.rur}；
\path{../Pgl50h/verification/pds_pc_control_write_ram_20261022/prj_tasks/pnr_1/place_route/hdmi_ddr3_top.ror}；
\path{chapter6/references/realtime_resource_audit_20261025.json}}
\endgroup

\subsection{模块资源归因与计算开销}
层次综合报告可以定位资源集中在哪些功能。
表~\ref{tab:rt-modules} 将较早实现与共享写缓冲实现按同名模块对照；
父模块用量包含子模块，表中各行不能相加作为全设计总量。

\begin{table}[htbp]
\centering\small
\caption{10 月 21 日至 22 日的模块资源变化}\label{tab:rt-modules}
\begin{tabularx}{\textwidth}{Yrrr}
\toprule
模块与统计项 & 10 月 21 日 & 10 月 22 日 & 变化 \\
\midrule
帧缓存整体 LUT / FF & 10,157 / 19,834 & 6766 / 7574 & 逻辑与寄存器减少 \\
分块渲染器 LUT / FF & 5202 / 10,149 & 5442 / 6076 & LUT 略增，FF 减少 \\
分块渲染器 DRM 等效量 & 101.5 & 83.5 & 减少 18 \\
投影器 LUT / FF & 4454 / 5439 & 4443 / 5440 & 基本不变 \\
投影器 APM & 43 & 43 & 不变 \\
DDR 仲裁器 LUT / FF & 4202 / 4684 & 350 / 592 & 减少 3852 / 4092 \\
DDR 控制器及 PHY LUT / FF & 6074 / 5076 & 6070 / 5076 & 基本不变 \\
\bottomrule
\end{tabularx}
\end{table}

投影器使用 43/45 个全设计 APM，占 95.56\%，说明算术资源主要服务于逐像素几何计算。
而源、坐标、输出块缓存和 FIFO 主导 DRM 用量，二者对应不同优化方向。
DDR 控制器及 PHY 在 10 月 22 日占 6070 LUT、5076 FF，
约为全设计综合 LUT 的 41.66\% 和 FF 的 33.83\%；
这些是支撑视频搬运的固定平台开销，不能全部归入矫正算法本身。

仲裁器 LUT 减少 91.67\%、FF 减少 87.36\%，投影器用量却几乎不变，
说明此轮资源收益主要来自数据暂存与传输组织的改进。
分块渲染器 LUT 从 5202 略增至 5442，也表明资源优化存在类型间交换，
并非所有模块、所有资源同时下降。
\begingroup\interlinepenalty=10000
\sourcefile{\path{../Pgl50h/prj_tasks/syn_1/synthesize/hdmi_ddr3_top.snr}；
\path{../Pgl50h/verification/pds_pc_control_write_ram_20261022/prj_tasks/syn_1/synthesize/hdmi_ddr3_top.snr}}
\endgroup

\subsection{片内缓存与片外帧存储}
完整图像置于 DDR，片内只保留正在处理的块和跨域数据。
主要缓存的逻辑数据容量见表~\ref{tab:rt-buffers}。

\begin{table}[htbp]
\centering\small
\caption{128$\times$32 配置的主要逻辑缓存容量}\label{tab:rt-buffers}
\begin{tabularx}{\textwidth}{Yp{5.1cm}r}
\toprule
缓存 & RTL 组织 & 容量 \\
\midrule
双槽源区域缓存 & $2\times2048\times256$ bit & 128 KiB \\
双槽坐标缓存 & $2\times4096\times23$ bit & 23 KiB \\
双槽输出块缓存 & $2\times256\times256$ bit & 16 KiB \\
采集及显示异步 FIFO & $2\times512\times256$ bit & 32 KiB \\
以上块 RAM 类缓存小计 & 不含小表及控制状态 & 199 KiB \\
采集及渲染写突发 RAM & $2\times16\times256$ bit & 1 KiB \\
\bottomrule
\end{tabularx}
\end{table}

源区域缓存占上述 199 KiB 的约 64.32\%，是主要存储项。
单个 64 KiB 源槽比 8 KiB 的最大输出块大得多，
因为源缓存必须容纳变换后的包围矩形、对齐填充和未直接采样的像素，
不能按输出块像素数直接缩减。
坐标缓存每像素存有效位及两维 11 bit 坐标，共 23 bit；
它以额外存储换取投影结果复用，使采样阶段无需再次执行完整逆投影。

器件全部 134 个 DRM18K 的标称容量约为 301.5 KiB，
远小于一帧 RGB565 的 4050 KiB。
199 KiB 只是几项 RTL 有效数据位的合计，并非综合推断的全部 RAM 容量。
实际映射还受宽深组合、端口模式、最小块粒度及其他查找表影响，
因此不能简单用 $199/2.25$ 估算 DRM 位置数。

DDR 中四个帧槽的有效像素载荷为
\begin{equation}
S_{\rm frames}=4S_f=16,588,800\ \text{bytes}=15.8203125\ \text{MiB}.
\end{equation}
每槽地址间隔为 16 MiB，四槽窗口为 64 MiB；
地址预留窗口与实际图像载荷并不相同。
RGB888 单帧需要 6,220,800 bytes，使用 RGB565 后，
相同像素访问次数下的存储和载荷均减少 33.33\%。
代价是红、蓝通道保留 5 bit，绿色保留 6 bit，输出位扩展不能恢复已丢失的信息。
因此这一资源收益应与第~\ref{chap:performance}~章的颜色与画质评价共同解释。

\section{优化收益及其性能代价}
\subsection{分块尺寸调整的存储收益}
从 $128\times64$ 改为 $128\times32$ 后，双槽坐标缓存从 46 KiB 降至
23 KiB，双槽输出缓存从 32 KiB 降至 16 KiB，合计减少 39 KiB 逻辑存储。
双槽源区域缓存仍保持 128 KiB，避免直接压缩姿态变换所需的源矩形容量。
原始层次报告中渲染器 DRM 等效量从 101.5 降至 83.5，与存储缩减方向一致；
但逻辑位数变化与物理位置变化仍应各按自身口径统计。

较小块增加了调度次数和边界访问，因此不能单独由存储收益推断速度改善。
参考资料记录早期 $128\times32$ 实现在一次压力测试中耗时 3,349,830 拍，
约 33.49830 ms，超过 30 fps 预算；10 月 22 日版本在改进同步 RAM 连续装载及
共享写数据通路后，完整并发测试最慢降至 33.14614 ms。
两次测试的 RTL 和覆盖条件不同，这说明小块必须与数据通路改造配合，
不能据此将时间差全部归因于块高变化。

\subsection{共享写缓冲与连续行装载}
旧式写通路复制整笔 $16\times256=4096$ bit burst，
并通过宽移位寄存器或选择网络推进数据。
新实现由采集端与渲染端各持有 $16\times256$ bit 分布式 RAM，
仲裁器仅保留当前所有者、beat 索引和事务状态，按索引读取当前客户端数据。
这减少重复宽数据暂存、复位扇出和数据选择互连，
与表~\ref{tab:rt-modules} 中仲裁器的大幅缩减相吻合。

写缓冲必须一直保持到最后一个有效数据传输完成，
只有 \texttt{WREADY} 有效才推进索引，并同时检查 \texttt{WLAST}。
等待期间数据稳定，完成后客户端才重填缓冲。
这一所有权规则使资源优化能够承受背压，避免仅在无停顿条件下正确。
大容量 RAM 不逐项复位，而通过状态、指针和先写后读约束保证有效性，
也有助于保持 RAM 推断。

结果行装载采用连续同步 RAM 流水：首拍发起读，此后每拍接收上一拍数据并发起下一次读，
排空最后一拍后提交写命令。其作用是减少逐字等待，提高已有 RAM 带宽利用，
而非增加几何计算并行度。共享写缓冲也没有改变投影公式和默认最近邻采样方式。

\subsection{物理资源的综合改善与扩展约束}
以 10 月 21 日原始物理报告和 10 月 22 日配套归档为同工具对照：
LUT 从 19,352 降为 16,878，减少 2474 个、约 12.78\%；
CLM FF 从 25,756 降为 14,871，减少 10,885 个、约 42.26\%；
DRM 位置从 122 降为 104，减少 18 个、约 14.75\%。
DRM 占用率下降 13.43 个百分点，可用位置由 12 增为 30。
这些收益来自分块、控制与写数据组织的整体变化，不能视为单因素消融结果。

APM 始终为 45，说明当前优化主要释放存储与互连资源。
39 个空余 APM 和 30 个空余 DRM 位置提供一定扩展空间，
但接入双线性采样、更大姿态范围或更高分辨率时，
还会同时改变缓存访问端口、源矩形容量、数据带宽与帧预算。
例如双线性需要周围采样点及插值计算，不能只按剩余 APM 数判断能否保留 30 fps。
目前没有功耗定量测量，亦不能由 FF 数下降直接给出能效提升百分比。

\section{静态时序与实现可靠性}
\subsection{时序余量与帧预算是两类约束}
100 MHz 渲染域每个时钟周期为 10 ns；静态时序分析检查寄存器路径能否在该周期内
满足 setup、hold 等要求，整帧分析检查多少个周期才能完成一幅图像。
通过其中一项不能替代另一项：帧周期数充足但 setup 违例的实现，
不能据 RTL 仿真确认其硬件运行可靠；所有路径满足时序也不保证处理周期数足够少。

工程 \path{prj_tasks/} 的 10 月 21 日 RTR 给出双角
\textsf{All Constraints Met}，慢角 DDR/渲染域 setup 余量 +0.462 ns；
10 月 22 日配套归档为 +0.873 ns；
10 月 24 日参考复核为 +1.245 ns。
这三个值属于不同实现，不能拼接成同一候选的时序结果。
10 月 24 日各项最小余量见表~\ref{tab:rt-timing}。

\begin{table}[htbp]
\centering\small
\caption{10 月 24 日 PDS2025.2 复核记录的最小时序余量（ns）}\label{tab:rt-timing}
\begin{tabularx}{\textwidth}{Yrr}
\toprule
检查类型 & Slow corner & Fast corner \\
\midrule
Setup & +0.721 & +0.121 \\
Hold & +0.114 & +0.053 \\
Recovery & +3.118 & +4.235 \\
Removal & +0.649 & +0.392 \\
Minimum pulse width & +0.397 & +0.568 \\
\bottomrule
\end{tabularx}
\end{table}

全部域最小 setup 为 +0.121 ns，并非渲染域的 +1.245 ns。
因此，不能只引用 100 MHz 域余量就认为所有接口均有相同裕量。
同样，0.187193 ms 的整帧余量与纳秒级路径余量分别约束调度和物理路径，
两者不能相互折算为统一的“系统裕量”。

\subsection{流水分级与 CDC 的配套检查}
参考报告中的历史优化将源矩形宽高计算、面积乘法和容量比较分为
\texttt{P\_BOUNDS}、\texttt{P\_AREA}、\texttt{P\_CAPACITY} 三个阶段，
用寄存器隔开长组合路径。对当时 $128\times64$ 的 255 个块，
每有效块增加两拍的名义总开销最多为 510 拍，即 5.10 $\mu$s，
约占 30 fps 帧预算的 0.0153\%。
若对 510 块配置作相同的逐块两拍预算，则为 1020 拍、10.20 $\mu$s，
这只是调度开销推算，不能替代该版本的实测增量。
这类分级以较少周期换取更短组合路径，其帧率代价通常小于取消关键流水级带来的时序风险。

异步 FIFO 使用 Gray 指针和双级同步器，但仍需约束各 Gray 位到首级同步器的偏斜。
10 月 24 日审计记录四条 10 bit Gray 总线均完成快、慢角检查，
最大偏斜 1.491 ns，小于 6.734 ns 约束。
10 月 22 日对应值为 1.470 ns，属于另一份归档。
\path{prj_tasks/} 的文本 RTR 中 Skew Summary 为空，
不能据此直接证明 Gray 总线位覆盖和偏斜全部合格，仍需独立审查偏斜数据库。
跨域功能仿真、同步结构、时序例外与偏斜检查分别覆盖不同问题，应共同作为依据。
\begingroup\interlinepenalty=10000
\sourcefile{\path{../Pgl50h/prj_tasks/pnr_1/report_timing/hdmi_ddr3_top.rtr}；
\path{chapter6/references/report.md}；
\path{chapter6/references/realtime_resource_audit_20261025.json}}
\endgroup

\subsection{跨工具重编译与未采用候选}
10 月 25 日审计记录：同一归档源码在 PDS2022.2-SP3 下重编译，
进程退出码为 0，也生成了位流，但慢角 DDR/渲染域 setup WNS 为
$-0.975$ ns、TNS 为 $-10.571$ ns，16 个端点失败。
该次综合/物理 LUT 为 13,280/13,321，数值虽低于 PDS2025.2 结果，
却伴随时序不通过，不能以更小 LUT 数宣称获得了可交付的资源优势。
不同工具版本和映射流程本身也会改变计数，二者不是单一 RTL 优化对照。

该次 FIFO 偏斜检查通过、最大偏斜为 1.724 ns，
说明 CDC 偏斜合格与同步逻辑 setup 合格是独立条件。
结构检查还列出 34 个无输入延迟和 86 个无输出延迟的条目，
需要按时钟、复位、异步控制与数据端口逐项解释，不能将进程正常结束或位流生成
解释为所有外部接口都获得了时序保证。

同一审计中另有针对帧缓存与渲染器的时序修复候选，记录为时序、CDC 和功能回归通过，
但明确标记 \texttt{not\_adopted}，主工程未采用，也未完成板测。
审计中的缩小分辨率功能回归亦不能代替后续源码的完整 1080p 吞吐重测。
因此，本章以 PDS2025.2 归档评价对应实现，以旧工具失败和未采用候选说明可复现边界，
不将不同快照的通过标记合并为当前版本的完整验收结论。

\section{性能结论与后续测量重点}
已有证据表明，完整帧存于 DDR、片内双槽分块复用和采集/渲染/显示重叠，
使归档 $128\times32$ 实现在指定 DDR 压力模型下达到约 30 个新结果帧/秒。
最慢观测渲染为 33.14614 ms、帧预算余量约 0.56\%，
说明实时目标已经在该次测试中达到，但进一步增加未被隐藏的等待空间有限。

资源方面，同工具物理对照显示 DRM 从 122 个位置降至 104 个，
FF 降幅约 42.26\%，收益主要来自块缓存缩减和共享写数据组织。
几何投影仍占绝大部分 APM，片内存储仍是主要数据资源限制；
后续优化应优先减少重复访存、缩短块间等待并改善服务间隔，
而不是单纯提高某条像素流水线的标称启动速率。

最终定量评价需要把实际采用的 RTL、约束、工具版本和位流哈希绑定，
在同一版本上补齐完整 1080p 多姿态压力回归、逐客户端 DDR 计数和板级帧身份测量。
应同时记录新结果帧率、最大渲染与发布间隔、端到端延迟分布、
FIFO 欠载/溢出及长时间运行情况。
这样才能将现有模型中的实时能力与资源收益，进一步落实为实际硬件的持续性能结论。
